\section{Conclusion}

We measured what a bounded, black-box text perturbation buys an attacker
against three small language models. These models work zero-shot for security
triage. Under a shared 15\% edit budget, random character-level perturbation
flips 0.0--5.0\% of the messages a model had already caught. Random
word-level perturbation flips 1.0--5.1\% while changing about half as many
units. Ordering the same word edit by measured saliency roughly doubles that:
4.6\% against 2.7\%, 33 message-specific flips against 7 when both are paired
on the messages they share.

Three things follow. First, a robustness number is a property of the search as
much as of the model. Inside the same 15\% word budget, choosing substitutions
rather than drawing them nearly doubles the flip rate (2.7\% to 4.6\%). This
result is not on the same messages or from the same candidate vocabulary. Once
we account for the higher budget actually spent, the factor is 1.4, not 1.7.
The search, the candidates, and the denominator must all be named in the claim.
Second, at this budget the binding constraint on a small triage model is what
it catches before any attack. What orders the flip rate is not accuracy and not
the measured margin: the model that flips most is neither the least accurate nor
the closest to its threshold, and its probability channel saturates, so the
margin cannot be read on it at all. The model that catches almost everything is
immovable because its decision rule is nearly constant. Detection quality, not
adversarial robustness, is where deployment effort pays first. Third, cheap
defences split cleanly. Normalising visually confusable characters removes every
character-level flip and only a small part of the word-level ones (7.5\% to
6.9\%). Re-thresholding to hold a false-positive budget can create evasions on
a model whose scores are least calibrated.

The measurement lesson runs through all of it. The transfer we observed was low
(0.9\% of caught messages off the diagonal). Per-model numbers are worth
measuring separately. The probability channel and the generated verdict agreed
on 941 of 1{,}001 rows. They diverged systematically on two of three models.
A single channel is not a measurement until it is checked against the other.
Every tool in this paper reports its denominator, its interval, and the bugs
we fixed in our own pipeline. Each bug silently inflated the attack success rate
in the direction that makes a paper look interesting.

The matched-vocabulary, matched-budget ablation of \S\ref{sec:ablation} is the
one experiment here that attributes rather than describes. With the candidate
set, the probe phase, and the call budget held fixed, it reports what ordering
is worth on its own. What it cannot do is turn five contrasts on roughly a
thousand paired messages into a ranking of attacks. \S\ref{sec:futurework}
names the two attacks that would test what remains; neither is reported here.
